\label{cap:revisao-bibliografica}

Na literatura acadêmica, observa-se uma lacuna significativa em estudos que aplicam técnicas de Otimização ao contexto de Alocação de Médicos em Clínicas de Exames. Esta carência é especialmente evidente nos trabalhos de \citeonline{Abdalkareem2021_ac} e \citeonline{ERHARD20181}, que, ao analisarem a literatura existente, destacam a falta de contribuições focadas especificamente nesse problema, enquanto o trabalho de \citeonline{Ngoo_Goh_Sze_Sabar_Abdullah_Kendall_2022} reforça a tendência contemporânea de estudos em problemas de Alocação de Enfermeiras (os).

O estudo de \citeonline{Abdalkareem2021_ac} apresenta uma análise abrangente de 190 trabalhos publicados entre 1995 e 2016, classificados de acordo com diferentes contextos. Entre as categorias destacam-se: o Problema de Agendamento para Admissão de Pacientes (PASP, do inglês \textit{Patient Admission Scheduling Problem}), o Problema de Alocação de Enfermeiras(os) (NRP, do inglês \textit{Nurse Rostering Problem}), o Agendamento de Salas de Operações (ORSP, do inglês \textit{Operation Room Scheduling Problem}), e uma categoria genérica denominada Outros, que abrange problemas como Alocação de Médicos(as), Escalas em Telemedicina e Assistência Médica Domiciliar.

Para cada um desses problemas, os autores analisam as metodologias utilizadas, que variam entre métodos exatos de programação matemática, técnicas heurísticas e meta-heurísticas. Além disso, examinam o tipo de base de dados empregado nas pesquisas, sejam elas baseadas em dados reais, simulações ou instâncias artificiais. O estudo revela que, na categoria ``Outros'', apenas três trabalhos tratam especificamente do problema de Alocação de Médicos(as), e todos eles focam em contextos hospitalares, empregando como principal metodologia a programação inteira mista.

Em um contexto mais focado, o trabalho de \citeonline{Ngoo_Goh_Sze_Sabar_Abdullah_Kendall_2022} busca analisar e avaliar a vasta literatura dos problemas de Alocação de Enfermeiras(os), focando em 50 publicações entre os anos de 2012 a 2021 com o intuito de entender as tendências e metodologias aplicadas para a solução dos problemas levantados. Além da definição formal do NRP e observações quanto a modelagem básica do problema, dada por sua extensibilidade e complexidade em restrições que buscam conciliar preferência com demanda e restrições de tempo. Também é analisada a tendência em métodos de solução, principalmente para programação linear inteira como as  mat-heurísticas (do inglês \textit{matheuristics}) aplicadas em contextos de estágio único (planejamento de horizontes de eventos determinísticos) e multiestágio (planejamento de horizonte de eventos estocásticos).

O trabalho de \citeonline{ERHARD20181} também visa categorizar e avaliar estudos, dessa vez aplicados ao contexto de alocação de médicos(as) por meio de métricas quantitativas. O artigo diferencia os trabalhos conforme o horizonte de planejamento: problemas de alocação estratégica, que tratam de decisões de longo prazo, como a definição e quantificação da equipe necessária; problemas de escalonamento, que atuam no nível tático, criando escalas de trabalho para o médio prazo; e problemas de replanejamento, que abordam ajustes de curto prazo, com foco na adaptação de escalas em função de demandas não previstas. A revisão de \citeonline{ERHARD20181} abrange 68 publicações entre 1985 e 2016.

O estudo também aponta que a modelagem matemática desses problemas, fora do escopo puramente financeiro, apresenta algumas características comuns. Dentre elas destaca-se o atendimento da demanda, a definição de limites inferiores e superiores na carga horária (referido no texto como corredor de carga horária), a garantia de descanso mínimo e o equilíbrio na distribuição da carga de trabalho entre a equipe médica. Além disso, os autores identificam a quantidade limitada de pesquisas focadas na área médica hospitalar, especialmente em comparação com os estudos sobre alocação de enfermeiras (os), uma tendência já observada no trabalho de \citeonline{Abdalkareem2021_ac}. Por fim, \citeonline{ERHARD20181} ressaltam algumas lacunas na literatura existente, como a falta de flexibilidade nas escalas e a pouca atenção dada à integração de turnos flexíveis e pausas, abrindo espaço para futuras investigações e aprimoramentos na área.

\section{Sobre a Literatura de Alocação de Médicos(as)}
\label{2rev-sec1}

Problemas de Alocação de Médicos(as) (do inglês \textit{Physician Allocation Problems}, abreviado como PAP) podem ser datados na literatura acadêmica desde o século passado, com trabalhos como o de \citeonline{Canter_1977} que realiza simulação para agendamento paciente-médico(a). Apesar disso, investigações voltadas a métodos de otimização, na maioria das vezes, enfatizavam contextos amplos, como a composição (\textit{rostering}) de equipes hospitalares, ou em alocações de enfermeiras (os), como o trabalho de \citeonline{Fitzpatrick_Farrell_Richter_Zeunik_1987}, com foco em suporte à tomada de decisão auxiliada por computador, com uma breve menção a possibilidade de implementação de métodos de otimização para solução da rotação de alocação de leitos.

\citeonline{Bruni_Detti_2014} avaliam o uso de modelos exatos por meio da implementação de um MIP (Programação Inteira Mista, do inglês \textit{Mixed Integer Programming}) no contexto da saúde, buscando formalizar uma solução generalista aplicável em hospitais, clínicas e outros. Embora a natureza de problemas de alocação de agentes sejam considerados NP-difícil \cite{Even_Itai_Shamir_1976}, o modelo proposto apresenta flexibilidade o suficiente para aplicação em diferentes contextos sem quebra de restrições obrigatórias (também referidas como \textit{hard constraints} na literatura) para sua solução mat-heurística baseada em \textit{branch-and-cut} para gerar soluções com melhoras significativas em relação à alocação manual. Testado como aplicação real no hospital universitário Umberto I, em Roma, a solução proposta melhorou consideravelmente a equidade nas alocações, reduzindo cargas de trabalho injustas (como múltiplas alocações do mesmo funcionário à suscetivos turnos noturnos) com eficiente boa eficiência computacional para aplicação em horizontes de planejamento extensos em poucos segundos.

A produção de \citeonline{Beaulieu_Ferland_Gendron_Michelon_2000} é uma das primeiras a tratar especificamente dos PAPs por meio de métodos de programação matemática e otimização. Ao conduzir um estudo de caso no hospital \textit{Sacré-Coeur}, a tarefa de gerenciar e escalar os horários dos 20 médicos(as) associados(as) se mostrou complexa devido a restrições obrigatórias complexas, tais como limites de turnos consecutivos, ergonomia de trabalho (como um período mínimo de descanso), períodos de férias, e restrições flexíveis, como preferências individuais por horários e dias de trabalho, embora a classificação das restrições dependa diretamente do hospital, dificultando a generalização do modelo. O trabalho propõe como metodologia o uso de programação linear inteira  multiobjetivo, resolvido através do método \textit{branch-and-bound}, buscando minimizar as violações ao cronograma, respeitando tanto as restrições rígidas quanto tentando atender as flexíveis.

Similarmente \citeonline{Topaloglu_2006} propõe um modelo exato baseado em \textit{goal programming} (método de programação linear multiobjetivo com foco em solucionar objetivos conflitantes) para solucionar o problema de escala em alas hospitalares emergenciais, focando especificamente na população médica residente, etapa necessária para obtenção do título de especialista em medicina. O estudo foi desenvolvido em parceria com um hospital universitário em sua EMR (ala de medicina residente emergencial, do ingles \textit{emergency medicine resident}). Devido sua natureza, o funcionamento da ala é ininterrupto, com grande fluxo de pacientes, gerando um desafio difícil para o controle e escalonamento de profissionais quando considerado regras ergonômicas de descanso, períodos de trabalho, considerações estudantis e preferências. Para gerir essas restrições foi implementado o uso de AHP (Processos Analíticos de Hierarquia, do inglês \textit{Analytical Hierarchy Process}) para gerir e priorizar o uso e atendimento de restrições.
Aplicado em uma base de dados de 15 meses, a abordagem proposta apresentou bons resultados em termos de equidade na escala de trabalho e conformidade quanto as restrições obrigatórias, principalmente quando comparados ao método manual na maior instância disponibilizada (5633 restrições e 23 residentes para um período de planejamento de 9 meses), embora o modelo não atingisse otimalidade nas especificações computacionais, especialmente tempo.

\citeonline{Beaulieu_Ferland_Gendron_Michelon_2000} e \citeonline{Topaloglu_2006}, quando comparadas ao escalonamento manual durante o período de implementação, tenham resultado em melhorias significativas na distribuição de horários e na redução do descumprimento de restrições, é importante destacar que o contexto do trabalho se aplica a salas de emergência, onde a demanda estática — isto é, demanda previamente conhecida e fixada no horizonte de planejamento — é irrelevante. Ou seja, o(a) médico(a) deve ser alocado(a) à sala independentemente de haver pouca demanda na ala emergencial, pois a chegada de pacientes é imprevisível.

Passando para um novo contexto, o estudo de \citeonline{Gunawan_Lau_2013} aborda o MPSP (Problema de Alocação do(a) Médico(a) Chefe, do inglês \textit{Master Physician Scheduling Problem}), que envolve a atribuição de médicos(as) a uma ampla gama de funções, como cirurgias, clínicas, tarefas administrativas e plantões, ao longo de um determinado horizonte de planejamento.
Para lidar com a abundante quantidade de tarefas a serem planejadas, os autores propuseram 4 modelos de programação linear inteira mista (MIP), sendo tais: \textit{modelo I} que visa minimizar tarefas não-alocadas, de maneira básica, desconsiderando pontos como preferência; \textit{modelos IIa e IIb}, que, respectivamente priorizam restrições flexíveis de preferência dos especialistas (chamado de agenda ideal) enquanto mantém o objetivo do primeiro modelo, e, procurar reduzir fadiga médica forçando restrições referentes a corredores de tempo de descanso e finais de semana (similar a definição de \citeonline{ERHARD20181}); enquanto o \textit{modelo III} procura um meio-termo entre as preferências dos modelos anteriores, por soluções ótimas de Pareto. Apesar de o uso de métodos heurísticos em instâncias maiores gerar soluções com desvios da otimalidade, os resultados eram melhores do que os gerados manualmente pela chefia de setor sem o auxílio de tomadores de decisão, concluindo a eficácia de MIPs na modelagem de problemas da área.

Os autores também observam que durante a execução do projeto, a restrição de disponibilidade de recursos estava diretamente atrelada a qualidade da solução, isto é, quando se limitava recursos nas restrições nos testes, aumentava-se o número de atividades não programadas. Apesar do que foi observado, a solução apresentada foi robusta o bastante para diversas instâncias.

\section{O Contexto Clínico}
\label{2rev-sec2}

Embora os esforços da comunidade cientifica, trabalhos como os de \citeonline{Beaulieu_Ferland_Gendron_Michelon_2000} e \citeonline{Topaloglu_2006}, assumem um contexto de sistema fechado, hospitalar, cujo intuito é garantir atendimento contínuo e ininterrupto ao paciente, respeitando questões interpessoais do profissional. Nesse contexto, um paciente já admitido pode ser atendido por vários profissionais durante sua estadia, remetendo a capacidade da estrutura hospitalar. Embora o estudo de \citeonline{Bruni_Detti_2014} seja generalizável para contextos da saúde, como o contexto ambulatorial (atendimentos rápidos e sem prévio agendamento), assume-se que o(a) especialista está associado(a) a uma única unidade (hospital). Nessa situação basta conhecer a demanda desta única facilidade.

Alguns autores, como \citeonline{ChenLinPeng2016} estendem as possibilidades do uso de modelagem matemática, aplicando otimização em dois estágios ao contexto hospitalar ambulatorial voltado ao funcionamento dos setores de exames clínicos, como radiologia e diagnóstico por imagem, na premissa de integrar os problemas de dimensionamento da equipe médica e alocação da equipe médica, enquanto \citeonline{BentayebLahrichiRousseau2023} estuda, de maneira similar, a integração de geração do quadro de consulta do paciente simultâneo à alocação de horários da equipe técnica em radiologia e exames por imagem para exames clínicos hospitalares. Em ambos os casos foram observadas melhorias significativas interpessoais da população alvo, como economia operacional devido à redução do dimensionamento da equipe necessária, além da aceleração no tempo de computação e confecção das escalas antecipadas. Somado a isso, a redução de violações de preferência, o trabalho de \citeonline{ChenLinPeng2016} apresenta desenvolvimento importante na literatura ao resolver o problema integrado, aumentando eficiência operacional e qualidade de vida, enquanto para \citeonline{BentayebLahrichiRousseau2023}, fora observada um aumento de 20\% no atendimento dos pacientes, reduzindo ociosidade de máquinas e encurtando custos operacionais correlatos. Os autores apontam, no entanto, possibilidades no futuro acadêmico de averiguar a robustez do sistema de dimensionamento de equipes, bem como a necessidade de generalização do estudo, que se limita ao contexto hospitalar médico. Também é amplamente discutido a hipótese de que a melhoria na eficiência de algoritmos também permitiria a expansão do estudo para toda a unidade, sem necessidade de segmentação, bem como a estabilidade da aplicação \cite{BentayebLahrichiRousseau2023}.


Para \citeonline{Edwards_Clague_Barlow_Clarke_Reed_Rada_1994}, foi necessário analisar a estrutura na qual os serviços se fundamentam de modo a identificar a origem dos problemas ambulatoriais de tempo de espera do paciente e fluxo de pacientes, métricas diretamente atreladas a satisfação do usuário final. Os autores averíguam por meio de simulações baseadas em dados reais os resultados obtidos pelo uso de sistemas de filas únicas - onde todos os pacientes são alocados a uma fila unitária - e filas quase-paralelas - onde são feitas triagem e criação de múltiplas filas de acordo com preferência e necessidade. Modelando tais sistemas como uma relação de agente (paciente) buscando por recurso limitado (médico(a) associado(a) a uma sala), foram observadas melhorias de 25\% no tempo de espera quando aplicadas filas quase-paralelas, devido sua eficiência e dinamismo, permitindo alterações e correções pontuais para o modelo, que por consequência aumentou a satisfação.
Em próprio texto íntegro, os autores comentam das lacunas no trabalho, especialmente na dificuldade de aplicação em contexto clínicos multidisciplinares (diferentes especializações) e a integração de um agenda de pacientes dinâmica, com reagendamentos, trocas de horários e faltas.

Outro trabalho que se dispôs similarmente a utilizar técnicas de simulação para alocar atendimentos em clínicas ambulatoriais foi \citeonline{Rau_Tsai_Liang_Tan_Syu_Jheng_Ciou_Jaw_2013}, exercido no contexto de sessões fisioterápicas, onde além das métricas de tempo de espera e satisfação, fora considerado a duração do exame por paciente, e consequentemente, a carga horária de trabalho do terapeuta. Com o uso de simulação de eventos discretos (do inglês, \textit{Discrete Event Simulation}, abreviada para DES), os autores obtiveram menores tempo de espera e um aumento nominal de 12 pacientes no fluxo diário, e especialmente à introdução de \textit{pooling} - na qual o paciente não é mais associado a um profissional preferencial, e pode ser atendido por outros profissionais da clínica - reduziu em 45\% os tempos de espera. Os autores \citeonline{Rau_Tsai_Liang_Tan_Syu_Jheng_Ciou_Jaw_2013} não consideram, no entanto, fatores estocásticos inerentes ao sistema, como falta de pacientes e desgaste dos equipamentos, limitando-se a estimativas baseadas em valores históricos.

Enquanto ambos os trabalhos citados não aplicam técnicas de otimização exata e modelo matemática para descriminar o problema observado, \citeonline{Alizadeh_Rezaeian_Abedi_Chiong_2020} aplicam MIP para construção de agendas de uso em exames não-emergenciais com necessidade de maquinário específico (referido na literatura como NOASP). Visando reduzir penalidades como: agendamentos em dias e horários indesejados, tempo de maquinário ocioso e trocas de médico(a), o modelo foca na preferência do paciente para criar agendas ótimas. Os autores apontam, no entanto, que o uso de métodos exatos se tornou inviável devido à natureza combinatorial do problema, quando escaladas as instâncias, levando-os a criar um algoritmo genético (abreviado para GA), que, embora não apresentasse soluções comprovadamente ótimas, se gerava agendamentos de boa qualidade, 10\% do desvio relativo.

\subsection{O Problema Multiclínicas}
\label{2rev-sec2-sub1}

Evidencia-se na Seção \ref{2rev-sec2} que  a maioria das investigações discorridas tem como principal motivação o estudo do agendamento em condições ambulatoriais (referido no inglês como \textit{outpatient}). Como não existem agendamentos prévios para esse sistema, trabalhos como o de \citeonline{ChenLinPeng2016} focam em equilibrar fatores estocásticos como incerteza de demanda por meio de otimização robusta, enquanto \citeonline{Alizadeh_Rezaeian_Abedi_Chiong_2020} consideram o(a) médico(a) como recurso a ser disponibilizado ao agente (paciente) para modelagem do sistema matemático; já os demais autores utilizam de artifícios de simulação com o(a) médico(a) fixo(a) para solução do problema. Os autores mencionados restringem-se a uma ótica unitária, onde há somente uma clínica, que, portanto, não é parte do modelo, mas sim delimitador do conjunto. Ao não considerar a possibilidade múltiplas unidades como parte do problema, abre-se mão da possibilidade de estudar o comportamento de modelos em situação onde um profissional pode ser alocado em diferentes ambientes, bem como suas implicações ergonômicas e legais.

O trabalho desenvolvido por \citeonline{Lai_Lu_Chen_2024} aborda o conceito de múltiplas clínicas em um hospital através do uso de \textit{frameworks} voltados a métodos de mineração de dados, como redes neurais artificiais e árvores de decisão para determinar mediante classificação os profissionais mais adequados a um grupo de pacientes. Como principal proposta do estudo, investigou-se a natureza de consultas que carecem de múltiplos segmentos, como encaminhamento para exames de ressonância magnética, os quais também compartilham recursos de outros setores do hospital, podem ser modeladas como grupos, que podem ser designados a combinação de equipes para melhor eficácia nos atendimentos.

\section{Heurísticas e Meta-Heurísticas em Escalonamento na Saúde}
\label{2rev-sec3}

Embora os métodos exatos de Programação Linear Inteira Mista discutidos nas seções anteriores ofereçam garantias de otimalidade, sua aplicabilidade é limitada quando a instância cresce em número de profissionais, períodos ou unidades, dada a natureza NP-difícil dos problemas de escalonamento \cite{Even_Itai_Shamir_1976}. Por essa razão, parcela significativa da literatura recente em alocação de médicos(as) e escalonamento de enfermeiras(os) recorre a heurísticas, meta-heurísticas e mat-heurísticas, abordagens que abrem mão da prova de otimalidade em troca de tempos computacionais viáveis em escala operacional, tendência também apontada nas revisões de \citeonline{Abdalkareem2021_ac} e \citeonline{Ngoo_Goh_Sze_Sabar_Abdullah_Kendall_2022}. Esta seção apresenta trabalhos primários representativos dessa linha, organizados em métodos heurísticos e mat-heurísticos, e na subseção subsequente em métodos evolutivos baseados em NSGA-II.

\citeonline{Vermuyten_Namorado_Rosa_Marques_Belien_Barbosa_Povoa_2018} abordam o escalonamento integrado de equipes em um serviço médico de emergência, considerando simultaneamente médicos(as), enfermeiros(as) e despachadores ao longo do horizonte de planejamento. A metodologia combina uma formulação em Programação Inteira com duas heurísticas baseadas em decomposição: \textit{diving} e Variable Neighborhood Decomposition Search (VNDS), aplicadas após relaxação parcial do modelo. Avaliada sobre instâncias do Instituto Nacional de Emergência Médica de Portugal, a heurística VNDS supera tanto o método \textit{diving} quanto um \textit{solver} comercial em qualidade de solução para horizontes realistas, e foi posteriormente integrada em sistema de apoio à decisão com interface gráfica para uso operacional. O trabalho evidencia a viabilidade de heurísticas estruturalmente alinhadas ao MIP em contextos onde a otimização exata se torna computacionalmente proibitiva.

\citeonline{Turhan_Bilgen_2020} propõem uma mat-heurística híbrida para o NRP, combinando os métodos Fix-and-Relax (F\&R), Fix-and-Optimize (F\&O) e Simulated Annealing (SA). O F\&R decompõe o horizonte em janelas semanais resolvidas como subproblemas MIP até a otimalidade, fornecendo soluções iniciais ao SA, que executa movimentos de vizinhança sobre escalas completas; o F\&O é invocado periodicamente para otimizar novamente subconjuntos de variáveis quando o SA estagna, alternando exploração global e refinamento exato local. Em testes sobre 24 instâncias \textit{benchmark} públicas, a abordagem híbrida supera os métodos do estado da arte na maioria dos casos e estabelece sete novos melhores resultados conhecidos. Os autores ilustram um arquétipo de mat-heurística no qual o MIP cuida da estrutura local enquanto o componente meta-heurístico responde pela diversificação, padrão metodológico diretamente relacionado à abordagem adotada neste estudo.

\citeonline{Ceschia_Di_Gaspero_Mazzaracchio_Policante_Schaerf_2023} introduzem uma formulação de NRP em hospitais italianos, considerando habilidades distintas, turnos especializados, limites de carga horária, tipos diversos de folgas, preferências individuais e incompatibilidades interpessoais. A solução combina um modelo MIP de referência com uma busca local guiada por SA que alterna dois movimentos: realocação enfermeiro-dia e troca de padrões de turno. Avaliada sobre 34 instâncias reais do norte da Itália, a metodologia produziu escalas aceitas pelas instituições e é distribuída como parte de \textit{software} comercial de \textit{rostering}. O trabalho exemplifica como uma meta-heurística orientada à prática absorve restrições idiossincráticas mantendo tempos de computação tratáveis para uso operacional cotidiano.

\citeonline{Santos_Toffolo_Gomes_Ribas_2016} apresentam técnicas de Programação Inteira para o NRP no contexto da \textit{International Nurse Rostering Competition} (INRC), combinando uma formulação compacta com geração de cortes e heurísticas primais baseadas em decomposição temporal. O método fixa as variáveis fora de uma janela deslizante de dias e reotimiza apenas o subproblema interno via MIP, padrão posteriormente reconhecido pela literatura como mat-heurística do tipo Fix-and-Optimize. A abordagem provou otimalidade na maioria das instâncias da competição, melhorou em até 15\% as melhores soluções conhecidas e venceu a INRC-2010 em todas as categorias. O trabalho consolidou a janela deslizante de F\&O como referência em escalonamento de enfermeiras(os).

\citeonline{Devesse_Santos_Toledo_2017} aplicam mat-heurísticas Relax-and-Fix e Fix-and-Optimize a um problema de escalonamento de médicos(as) em prontos-socorros, com formulação MIP \textit{minimax} orientada à equidade na distribuição de turnos indesejáveis sob restrições da legislação trabalhista brasileira. As heurísticas particionam as variáveis binárias por semanas e por tipos de turno, resolvendo cada subproblema com CPLEX e aceitando melhorias incrementais sobre a escala incumbente. Avaliada em instâncias com 10 a 40 médicos(as) e horizontes de 4 a 36 semanas, a abordagem mostra que R\&F e F\&O produzem soluções competitivas com o MIP em uma fração do tempo, embora o MIP completo ainda dominasse instâncias menores. O trabalho é particularmente relevante por aplicar diretamente o par R\&F + F\&O ao escalonamento médico, contexto idêntico ao tratado no Capítulo \ref{cap:mip-heuristicas}.

\citeonline{Wickert_Kummer_Neto_Boniatti_Buriol_2021} propõem um modelo de Programação Inteira e uma mat-heurística Fix-and-Optimize para o escalonamento de médicos(as) no Hospital de Clínicas de Porto Alegre, considerando preferências individuais, restrições legais, requisitos hospitalares e equilíbrio de carga de trabalho. O método constrói uma solução factível via MIP e iterativamente reotimiza vizinhanças definidas por dia, médico(a), semana ou turno, mantendo as demais variáveis fixadas. Em instâncias sintéticas com 50, 100 e 150 médicos(as) ao longo de 4 semanas, a mat-heurística F\&O sistematicamente supera o MIP em horizontes de 10 a 30 minutos, particularmente nas instâncias maiores em que o solver enfrenta dificuldade de convergência. O estudo documenta um precedente direto, em contexto brasileiro, do uso de F\&O para escalonamento de médicos(as) em larga escala.

Os trabalhos apresentam padrão recorrente, evidenciado em escalonamento de enfermeiras(os) e médicos(as): a estrutura combinatorial do MIP é preservada em subproblemas locais, enquanto componentes meta-heurísticos respondem pela diversificação no espaço de busca. Essa hibridização é particularmente útil em horizontes longos e em contextos com múltiplas restrições conflitantes, condições que também caracterizam o problema vigente, motivando a adoção de métodos análogos no Capítulo \ref{cap:mip-heuristicas}.

\subsection{Abordagens Multiobjetivo em Escalonamento}
\label{2rev-sec3-sub1}

Problemas de escalonamento na saúde combinam objetivos tipicamente conflitantes, como cobertura da demanda, equidade na distribuição de turnos e atendimento a preferências individuais, motivando o tratamento explícito da estrutura multiobjetivo da formulação. A literatura organiza-se em duas grandes famílias metodológicas: o escalonamento, que agrega múltiplos critérios em uma única função objetivo via custos ponderados ou programação por metas, e a busca direta de soluções Pareto-eficientes, suportada tanto por métodos exatos de decomposição quanto por meta-heurísticas baseadas em dominância. As subseções a seguir percorrem trabalhos representativos dessas duas famílias, culminando no NSGA-II, algoritmo evolutivo adotado como base no Capítulo \ref{cap:mip-heuristicas}.

\citeonline{Parr_Thompson_2007} estudam o NRP com função de custo ponderada, tendo a cobertura como restrição rígida e os demais critérios (preferências, padrões de turno, equidade) agregados em penalidades soft. A metodologia compara duas estratégias de busca local sobre a função ponderada: SAWing, que ajusta dinamicamente os pesos durante o \textit{simulated annealing}, e Noising-SA, que perturba os custos com ruído para escapar de ótimos locais. O trabalho demonstra que Noising-SA produz escalas com penalidade total consistentemente menor, sendo menos sensível à calibração de pesos, evidenciando uma das principais dificuldades práticas do escalonamento.

\citeonline{Maenhout_Vanhoucke_2010} formulam o NRP de médio prazo com múltiplos objetivos e prioridades distintas, modelando-o como um problema de \textit{set partitioning} resolvido por \textit{branch-and-price} exato. A contribuição central está nas estratégias de \textit{branching} e \textit{pruning} adaptadas ao contexto multiobjetivo, capazes de preservar a eficiência da geração de colunas enquanto direcionam a busca para escalas de alta qualidade. Os experimentos computacionais demonstram que regras de \textit{branching} bem projetadas viabilizam a otimização exata de instâncias realistas, oferecendo um contraponto exato à dominância de meta-heurísticas no escalonamento multiobjetivo.

\citeonline{Burke_Li_Qu_2012} propõem busca Pareto-direta para o NRP em uma UTI holandesa, com formulação que possui nove restrições suaves distintas. O método combina geração aleatória de padrões legais por enfermeiro, heurística adaptativa para alocação inicial, reparo de cobertura e \textit{simulated annealing} multiobjetivo com dois modos de avaliação: soma ponderada e dominância de Pareto. A variante baseada em dominância produz dezenas a milhares de escalas não dominadas em cada instância, em cenário onde a Programação Inteira completa não encontrou solução factível, ilustrando a vantagem da busca Pareto-direta em formulações fortemente restritas.

A linha evolutiva multiobjetivo passou a dominar a literatura subsequente em escalonamento na saúde a partir do NSGA-II, formalizado por \citeonline{Deb_Pratap_Agarwal_Meyarivan_2002}, cuja contribuição metodológica é detalhada a seguir.

O algoritmo NSGA-II foi proposto por \citeonline{Deb_Pratap_Agarwal_Meyarivan_2002} para superar três limitações dos algoritmos evolutivos multiobjetivo anteriores: complexidade computacional elevada, ausência de elitismo e dependência do parâmetro de \textit{sharing}. A metodologia introduz a ordenação não-dominada rápida com complexidade \(O(MN^{2})\), a distância de \textit{crowding} para preservação de diversidade, e a seleção elitista baseada na união entre populações pai e filha, além de uma regra de dominância modificada para problemas com restrições. Em \textit{benchmarks} de teste, o NSGA-II superou os algoritmos SPEA e PAES em convergência e em distribuição da fronteira aproximada de Pareto. O trabalho consolidou o NSGA-II como o algoritmo evolutivo multiobjetivo de referência em Pesquisa Operacional, sendo adotado ou utilizado como linha de base em grande parte da literatura subsequente em escalonamento.

\citeonline{Hamid_Tavakkoli_Moghaddam_Golpaygani_Vahedi_Nouri_2020} formulam um problema multiobjetivo de escalonamento de enfermeiras(os) que enfatiza fatores humanos, com três objetivos: minimização do custo total de pessoal, minimização da incompatibilidade entre estilos de tomada de decisão na mesma escala e maximização da satisfação com os turnos atribuídos. Para resolução, comparam três meta-heurísticas multiobjetivo - NSGA-II, algoritmo Keshtel multiobjetivo e Busca Tabu multiobjetivo - e utilizam Análise Envoltória de Dados (DEA) para classificar as soluções da fronteira de Pareto. Em estudo de caso em hospital iraniano, todos os algoritmos geram escalas factíveis com diferentes compromissos entre os critérios, sendo o NSGA-II competitivo em termos de qualidade da fronteira. Os autores evidenciam que tratar fatores humanos como objetivos explícitos altera estruturalmente o conjunto Pareto-eficiente, com pequenos aumentos de custo gerando ganhos significativos em compatibilidade e satisfação.

\citeonline{Lavygina_Welsh_Crispin_2019} aplicam o NSGA-II ao escalonamento de médicos(as) residentes no Reino Unido, sob regras nacionais de horas trabalhadas e descanso. A novidade do estudo é tratar a equidade como segundo objetivo, formalizada como o desvio padrão das penalidades acumuladas por médico(a), de modo que soluções equitativas distribuem características indesejáveis da escala uniformemente entre os profissionais. Comparado contra variantes de Pareto Simulated Annealing e algoritmos genéticos mono-objetivo em dez instâncias, o NSGA-II produziu fronteiras com hipervolume consistentemente superior e convergência mais rápida. A aplicação direta a médicos(as) e a inclusão da equidade como objetivo explícito o aproximam metodologicamente da estrutura de objetivos conflitantes adotada no problema vigente.

\citeonline{Khalili_Shahrezaei_Abri_2020} propõem um modelo bi-objetivo de escalonamento de enfermeiras(os) que incorpora o fator fadiga, modelado em função dos intervalos de descanso e dos padrões de turnos. Para a maior instância do estudo de caso (90 profissionais no hospital iraniano Labbafinejad), o modelo é resolvido pelo NSGA-II e por MOPSO, com cromossomos especificamente projetados, e o método TOPSIS combinado a ponderação por entropia é empregado para selecionar a solução final na fronteira de Pareto. O NSGA-II encontrou escalas de melhor qualidade em menor tempo computacional, com soluções que indicam padrões de descanso noturno e intermediário capazes de reduzir fadiga sem comprometer significativamente a estrutura de custos. O trabalho ilustra a combinação NSGA-II com técnicas de seleção a posteriori, padrão útil quando o tomador de decisão precisa de uma única recomendação a partir do conjunto de Pareto.

O padrão recorrente nas aplicações evolutivas multiobjetivo aqui descritas envolve representação direta das variáveis de alocação no cromossomo, operadores de cruzamento e mutação adaptados ao contexto e, frequentemente, a integração com técnicas auxiliares de seleção ou refinamento local. Esses elementos motivam diretamente as variantes propostas no Capítulo \ref{cap:mip-heuristicas}, nas quais o NSGA-II base é estendido com operadores customizados e adaptações estruturais para o problema de alocação de médicos(as) em múltiplas clínicas.

\section{Lacunas de Pesquisa}

Não é comum observar na literatura acadêmica investigações acerca da natureza de múltiplas unidades de um mesmo segmento clínico cujas operações sejam completamente autônomas entre si, como filiais de redes laboratoriais de exames. Nesse contexto, a unidade tem recursos (salas, equipamentos, médicos), disponibilidades e demanda intrínsecos a si mesma, e a relação de \textit{vários-para-vários} (como exemplificado na imagem da Figura \ref{fig:rel-n-to-n}) entre profissional e unidade causa uma explosão combinatória nas possibilidades de escalonamento, uma vez que existem múltiplas possibilidades de alocações possíveis, o que inevitavelmente gera situações de remanejamento, onde um profissional outrora alocado em uma unidade é escalado para outra. Além disso, as demandas e disponibilidades, bem como fluxo operacional, podem não ser iguais, dificultando a modelagem quando aplicada a um panorama de encontrar boas soluções para todas as clínicas de uma só vez.

Não é considerado também o contexto de especialistas associados em regime de trabalho de pessoa jurídica. Nessa situação as normas ergonômicas de trabalho e questões salariais são completamente diferentes do normalmente observado na literatura, assim como não foram constatados estudos acerca do remanejamento dos profissionais em relação a diversas unidades ao longo do horizonte de planejamento e suas implicações quanto a qualidade da solução.

\newpage

\begin{figure}[h]

	\begin{subfigure}{0.5\textwidth}
		\includegraphics[scale=0.1]{images/clinic-one-to-m.png}
		\caption{Relação vários para um}
		\label{fig:rel-n-to-1}
	\end{subfigure}
	\begin{subfigure}{0.5\textwidth}
		\includegraphics[scale=0.1]{images/clinic-many-to-many.png}
		\caption{Relação vários para vários}
		\label{fig:rel-n-to-n}
	\end{subfigure}

	\caption{Exemplificação da diferença entre os contextos de relacionamento apresentados.}
	\label{fig:exemplo-relacionamentos}
	\fautor
\end{figure}

Para este estudo de mestrado considera-se como problema a alocação de médicos de um único segmento clínico em múltiplas unidades, considerando que tanto o profissional com sua agenda de disponibilidade para atendimentos, e a clínica com sua demanda de exames a cumprir devem ser respeitadas, não sendo possível, portanto realizar a alocação em períodos onde a unidade ou o profissional não estão disponíveis. Aplica-se também a situação na qual a folga de um profissional deve ser coberta. Além disso, propõe-se a investigação acerca da possibilidade de redução nas trocas de clínicas trabalhadas por um profissional ao longo do horizonte de planejamento e suas implicações quanto ao atendimento de demanda.

Ao avaliar o sistema na totalidade, buscando atender critérios operacionais de maneira global, para todas as unidades simultaneamente, têm-se um sistema complexo de alta sensibilidade, onde a alocação de um profissional (agente) em uma clínica pode afetar outra, assemelhando-se a definição de escalonamento (do inglês \textit{rostering}) proposta por \cite{Wren_1996}.

O Quadro \ref{tab:taxonomia-estudos} de avaliação taxonômica ressalta de maneira sucinta aspectos-chave dos principais trabalhos a serem considerados na arguição deste estudo quanto às lacunas de pesquisa e descreve de maneira sucinta as contribuições deste estudo. Nele, o campo \textit{Escalas Acíclicas} indica se o trabalho considera escalas em que os profissionais não seguem um padrão fixo e repetitivo de turnos ao longo do horizonte de planejamento — característica presente no contexto estudado, onde cada médico(a) possui disponibilidades individuais e variáveis entre os dias.


\begin{quadro}[h]
	\centering
	\resizebox{\textwidth}{!}{%
		\begin{footnotesize}
			\begin{tabular}{|M{2cm}|M{2cm}|M{1.5cm}|M{1.5cm}|M{1.5cm}|M{1.5cm}|M{1.8cm}|M{3cm}| }
				\hline
				\textbf{Estudo}                                                        & \textbf{Ambientação} & \textbf{Múltiplas Unidades}    & \textbf{Alocação de Médicos} &
				\textbf{Escalas Acíclicas}                                             & \textbf{Dados Reais} & \textbf{Minimização de Trocas} & \textbf{Metodologia}                                                                                                               \\
				\hline
				\citeonline{Edwards_Clague_Barlow_Clarke_Reed_Rada_1994}               & Ambulatório          &                                &                              &            & \textbf{X} &            & Simulação                                                    \\
				\hline
				\citeonline{BentayebLahrichiRousseau2023}                              & \textbf{Exames}  &                                & \textbf{X}                   & \textbf{X} &            &            & \textbf{Modelagem Matemática}                                \\
				\hline
				\citeonline{ChenLinPeng2016}                                           & Ambulatório          &                                & \textbf{X}                   &            &            &            & \textbf{Modelagem Matemática}                                \\
				\hline
				\citeonline{Bruni_Detti_2014}                                          & Hospitalar           &                                & \textbf{X}                   &            &            &            & \textbf{Modelagem Matemática}                                \\
				\hline
				\citeonline{Gunawan_Lau_2013}                                          & Hospitalar           &                                & \textbf{X}                   &            &            &            & Heurística                                                   \\
				\hline
				\citeonline{Topaloglu_2006}                                            & Emergencial          &                                &                              & \textbf{X} & \textbf{X} &            & \textbf{Modelagem Matemática}                                \\
				\hline
				\citeonline{Alizadeh_Rezaeian_Abedi_Chiong_2020}                       & Ambulatório          &                                & \textbf{X}                   &            &            &            & Heurística                                                   \\
				\hline
				\citeonline{Rau_Tsai_Liang_Tan_Syu_Jheng_Ciou_Jaw_2013}                & Ambulatório          &                                &                              &            & \textbf{X} &            & Simulação                                                    \\
				\hline
				\citeonline{Lai_Lu_Chen_2024}                                          & Hospitalar           & \textbf{X}                     &                              &            &            &            & Mineração de Dados                                           \\
				\hline
				\citeonline{Vermuyten_Namorado_Rosa_Marques_Belien_Barbosa_Povoa_2018} & Emergencial          &                                & \textbf{X}                   &            & \textbf{X} &            & Heurística                                                   \\
				\hline
				\citeonline{Turhan_Bilgen_2020}                                        & Hospitalar           &                                &                              &            &            &            & mat-heurística                                                \\
				\hline
				\citeonline{Lavygina_Welsh_Crispin_2019}                               & Hospitalar           &                                & \textbf{X}                   &            & \textbf{X} &            & Meta-heurística                                               \\
				\hline
				\citeonline{Hamid_Tavakkoli_Moghaddam_Golpaygani_Vahedi_Nouri_2020}    & Hospitalar           &                                &                              &            & \textbf{X} &            & Meta-heurística                                               \\
				\hline
				\citeonline{Khalili_Shahrezaei_Abri_2020}                              & Hospitalar           &                                &                              &            & \textbf{X} &            & Meta-heurística                                               \\
				\hline
				\citeonline{Ceschia_Di_Gaspero_Mazzaracchio_Policante_Schaerf_2023}    & Hospitalar           &                                &                              &            & \textbf{X} &            & Meta-heurística                                               \\
				\hline
				\citeonline{Parr_Thompson_2007}                                        & Hospitalar           &                                &                              &            &            &            & Meta-heurística                                               \\
				\hline
				\citeonline{Maenhout_Vanhoucke_2010}                                   & Hospitalar           &                                &                              &            &            &            & \textbf{Modelagem Matemática}                                \\
				\hline
				\citeonline{Burke_Li_Qu_2012}                                          & Hospitalar           &                                &                              &            & \textbf{X} &            & Meta-heurística                                               \\
				\hline
				\citeonline{Santos_Toffolo_Gomes_Ribas_2016}                           & Hospitalar           &                                &                              &            &            &            & mat-heurística                                                \\
				\hline
				\citeonline{Devesse_Santos_Toledo_2017}                                & Emergencial          &                                & \textbf{X}                   &            &            &            & mat-heurística                                                \\
				\hline
				\citeonline{Wickert_Kummer_Neto_Boniatti_Buriol_2021}                  & Hospitalar           &                                & \textbf{X}                   &            & \textbf{X} &            & mat-heurística                                                \\
				\hline
				\textbf{Este estudo}                                                   & \textbf{Exames}   & \textbf{X}                     & \textbf{X}                   & \textbf{X} & \textbf{X} & \textbf{X} & \textbf{Modelagem Matemática, Heurísticas e Meta-Heurísticas} \\
				\hline
			\end{tabular}
		\end{footnotesize}}
	\caption{Análise taxonômica de estudos em Pesquisa Operacional relacionados à Alocação e Escalonamento na saúde.}
	\fautor
	\label{tab:taxonomia-estudos}
\end{quadro}